\begin{casebox}
\textbf{知识点小案例：CAS 成功不代表旧节点可以马上释放。}
Treiber stack 的特例是线程 A 读到 top 指向节点 X，准备 CAS；线程 B pop 出 X 并释放，又把同一地址复用成新节点；线程 A 的 CAS 可能看到“地址没变”而误以为状态没变。ABA 和回收问题说明：线性化点只解决状态更新顺序，不解决对象生命周期。由此推出规律：无锁证明必须同时写线性化点、ABA 防护和回收策略。工程上没有 hazard pointer、epoch、RCU 或 tagged pointer 这类方案时，不能把含指针回收的无锁结构放进默认路径。
\end{casebox}
